\chapter{Metodologia}
\label{ch:metodologia}

Este capítulo detalha os procedimentos metodológicos e o delineamento experimental adotados para alcançar os objetivos propostos neste trabalho de conclusão de curso. A pesquisa classifica-se como aplicada e quantitativa, visando a implementação, simulação e validação de um modelo de anonimização em dados de saúde.

A abordagem metodológica fundamenta-se no conflito teórico, explorado na fundamentação, entre a necessidade de proteção de dados sensíveis — imposta por legislações como a Lei Geral de Proteção de Dados Pessoais (LGPD) — e a crescente demanda por dados de alta fidelidade para o treinamento de algoritmos de Inteligência Artificial e apoio à decisão clínica. Para investigar esse fenômeno, adota-se um fluxo de trabalho experimental iterativo, que parte da seleção criteriosa de dados reais, passa pela aplicação controlada de algoritmos de privacidade e culmina na análise comparativa de performance de modelos preditivos.

O método proposto não se limita apenas à execução técnica da anonimização, mas busca estabelecer um protocolo de validação que possa ser replicado em outros contextos de pesquisa médica. As seções a seguir descrevem cada etapa deste processo, justificando as escolhas de ferramentas, dados e métricas com base na literatura científica atual.

\section{Delineamento Experimental}
\label{sec:metod-delineamento}

O experimento foi desenhado para testar a hipótese central de que é possível aplicar técnicas de anonimização robustas em dados clínicos complexos mantendo a viabilidade para o treinamento de modelos de \textit{Machine Learning} (ML). Diferentemente de abordagens que avaliam a utilidade dos dados apenas através de métricas estatísticas genéricas, este trabalho propõe uma validação orientada à tarefa (\textit{task-oriented}), onde a qualidade do dado é medida pela sua capacidade de resolver um problema clínico real.

Para estruturar essa investigação, o estudo é organizado em torno de variáveis controladas e observadas:

\begin{itemize}
    \item \textbf{Variáveis Independentes (Controle):} Correspondem aos parâmetros dos algoritmos de anonimização aplicados. Isso inclui a escolha da técnica (ex: k-anonimato, l-diversidade, Privacidade Diferencial), a intensidade da proteção (o valor de $k$, o valor de $\epsilon$) e as hierarquias de generalização definidas para cada atributo.
    \item \textbf{Variáveis Dependentes (Observadas):}
    \begin{enumerate}
        \item O \textbf{Risco de Reidentificação Residual}, medido pelas métricas de privacidade discutidas na \autoref{subsec:fund-metricas-privacidade}, que indicam a probabilidade de um ataque bem-sucedido.
        \item A \textbf{Utilidade dos Dados}, mensurada duplamente: pela preservação das propriedades estatísticas do conjunto (utilidade geral) e, principalmente, pela performance do modelo preditivo treinado (utilidade específica).
    \end{enumerate}
\end{itemize}

O experimento seguirá o método comparativo rigoroso. Será estabelecido um "Grupo de Controle", composto pelos dados originais apenas desidentificados (sem generalização), que servirá como a linha de base (\textit{baseline}) de performance ideal. Em contrapartida, serão criados múltiplos "Grupos Experimentais", compostos pelos mesmos dados submetidos a diferentes configurações de anonimização. A comparação direta entre o controle e os grupos experimentais permitirá isolar o impacto da anonimização, quantificando a degradação da performance do modelo de IA como uma função do aumento da privacidade.

\section{Aquisição e Definição da Amostra (MIMIC-III)}
\label{sec:metod-aquisicao}

% TODO: explicar porque usamos o MIMIC-III usando literatura e comparações

A validade de um estudo sobre anonimização depende intrinsecamente da qualidade, da granularidade e da complexidade dos dados utilizados. Para garantir a relevância prática dos resultados, este trabalho utilizou dados obtidos através da plataforma \textit{PhysioNet} \cite{Goldberger2000}, um repositório de referência mundial para pesquisa médica gerido pelo MIT Laboratory for Computational Physiology. O acesso foi obtido através de processo de credenciamento rigoroso que incluiu treinamento certificado em ética em pesquisa com seres humanos.

Foi selecionado o banco de dados \textbf{MIMIC-III v1.4} \cite{Johnson2016}, um repositório relacional contendo informações desidentificadas de aproximadamente 60.000 admissões em Unidades de Terapia Intensiva do \textit{Beth Israel Deaconess Medical Center}. A estrutura complexa, com tabelas interligadas de dados demográficos, sinais vitais, procedimentos e diagnósticos, oferece o cenário ideal para testar a robustez dos algoritmos de anonimização em contextos clínicos realistas.

% TODO: ter uma seção dedicada somente a limpeza de dados, entrando em detalhes e citando o preparo de dados do Harutyunyan2019

O processo de filtragem da amostra ocorreu em duas etapas. Inicialmente, restringiu-se a amostra a pacientes adultos ($\geq 18$ anos), estabelecendo um universo primário de $N_{\text{total}} = 42.277$ admissões. Esta base completa foi utilizada para os cálculos de proteção de privacidade na ferramenta ARX, visando fornecer aos algoritmos máxima diversidade demográfica para a auditoria de risco de reidentificação.

Em uma segunda etapa, foi aplicado o critério de duração mínima de permanência em UTI. O modelo ConCare requer séries temporais fisiológicas completas de 48 horas para seu funcionamento. Admissões com duração inferior a 48 horas foram removidas, resultando em uma amostra analítica final de $N_{\text{task}} = 17.917$ admissões para treinamento, validação e teste do modelo preditivo. A \autoref{tab:amostra-caracterizacao} apresenta a caracterização demográfica e clínica desta coorte final.

% TODO: checar se esses dados estão corretos

\begin{table}[ht]
\centering
\caption{Caracterização da Amostra Analítica Final ($N = 17.917$)}
\label{tab:amostra-caracterizacao}
\begin{tabular}{ll}
\toprule
\textbf{Característica} & \textbf{Valor / Estatística} \\ \midrule
Total de Admissões & 17.917 \\
Duração média de permanência (horas) & $73.4 \pm 52.1$ \\
Idade média (anos) & $63.9 \pm 17.5$ \\
Gênero (Female / Male) & 44.1\% / 55.9\% \\
Mortalidade hospitalar & 10.6\% (1.899 óbitos) \\
\midrule
\textbf{Principais Comorbidades (CCS)} & \textbf{Prevalência} \\
Essential Hypertension & 41.6\% \\
Cardiac Arrhythmias & 37.5\% \\
Heart Failure & 32.9\% \\
Coronary Atherosclerosis & 32.8\% \\
Fluid Disorders & 31.3\% \\
\bottomrule
\end{tabular}
\fonte{Elaborado pelo autor, conforme dados do MIMIC-III filtrado segundo critérios de inclusão/exclusão do estudo.}
\end{table}

\section{Modelo de Referência para Avaliação de Utilidade Clínica}
\label{sec:metod-concare}

Uma das maiores dificuldades na literatura de anonimização é definir o que constitui "utilidade aceitável". Métricas puramente estatísticas, como a "perda de informação" (\textit{Information Loss}), são abstratas e muitas vezes não refletem a viabilidade prática do dado para aplicações clínicas reais. Para resolver isso, este trabalho adota uma abordagem orientada à tarefa (\textit{task-based evaluation}). A premissa é simples: o dado é útil se ele permite resolver o problema clínico para o qual foi coletado.

% TODO: Justificar a escolha do modelo usando literatura e comparação

Foi selecionado o modelo \textbf{ConCare} \cite{Ma2020} como referência para avaliação de utilidade. O ConCare é uma arquitetura de rede neural que se distingue pela integração explícita de dados estáticos (demográficos) com séries temporais fisiológicas através de um mecanismo de \textit{Cross-Attention}. Esta escolha é particularmente rigorosa para testar o impacto da anonimização, pois os dados demográficos (Age, Gender, Ethnicity) — precisamente os atributos atacados pela anonimização — são componentes fundamentais do processamento do modelo. Uma degradação na precisão desses atributos impacta diretamente a capacidade do ConCare de ponderar a importância de momentos específicos durante a trajetória clínica.

% TODO: explicar porque predição de mortalidade foi a tarefa escolhida usando literatura e dados reais

A tarefa clínica selecionada é a \textbf{predição de mortalidade hospitalar}, um problema crítico de prognóstico em UTI, clinicamente relevante, e que depende essencialmente tanto de dados temporais (sinais vitais, exames laboratoriais) quanto de contexto demográfico estático. O desempenho do modelo ConCare treinado nos dados originais sem anonimização adicional estabelece a \textit{baseline} (linha de base) de performance ideal. As comparações posteriores com os dados anonimizados permitem quantificar a degradação de utilidade como uma função direta do aumento de proteção de privacidade.

\section{Definição de Quasi-Identificadores, Atributos Sensíveis e Outcome}
\label{sec:metod-categorias}

Para a aplicação eficaz de técnicas de anonimização, foi realizada a classificação semântica dos atributos conforme o modelo de análise de risco proposto por El Emam e Arbuckle \cite{ElEmam2011}. Na avaliação da privacidade, foram definidos os seguintes componentes:

\begin{itemize}
    \item \textbf{Quasi-Identificadores (QIs):} \textbf{Age} (idade em anos), \textbf{Gender} (sexo: Male/Female), e \textbf{Ethnicity} (etnia: categorias disponíveis no MIMIC-III). Estes três atributos foram selecionados como os alvos principais das técnicas de anonimização, pois sua combinação pode levar à reidentificação quando ligados a bases de dados externas.
    \item \textbf{Atributo Sensível (AS):} \textbf{Primary Diagnosis} (diagnóstico principal, codificado segundo ICD-9). O objetivo da anonimização é quebrar a associação determinística entre os QIs e este atributo, impedindo que um adversário infira o diagnóstico a partir de informações demográficas.
    \item \textbf{Outcome:} \textbf{Mortality} (resultado: óbito hospitalar sim/não). Este é o rótulo de classe utilizado para treinamento e avaliação do modelo ConCare.
    \item \textbf{Atributos Preservados:} \textbf{Gender} foi preservado em todos os cenários de anonimização, como consta na metodologia de proteção clínica, pois sua perda teria impacto excessivo na utilidade clínica sem ganho significativo de privacidade.
\end{itemize}

\section{Aplicação de Técnicas de Anonimização}
\label{sec:metod-anonimizacao}

% TODO: expandir explicação e jusiticativa para o uso da ferramenta, usando literatura e comparação com outras ferramentas

O processo de anonimização foi executado utilizando o \textbf{ARX Data Anonymization Tool} \cite{Prasser2014}, uma plataforma especializadamente projetada para anonimização de dados biomédicos. O ARX oferece implementações robustas de k-anonimato, l-diversidade, e t-closeness, além de métricas de risco integradas que facilitam a auditoria de proteção.

\subsection{Definição de Hierarquias de Generalização}

% TODO: expandir justificativa para essas definições usando dados reais e literatura

Para cada quasi-identificador, foram definidas hierarquias de generalização (taxonomias) que permitem a progressiva transformação dos valores, do mais específico ao mais generalizado, preservando significado clínico. As hierarquias foram configuradas da seguinte forma:

\begin{itemize}
    \item \textbf{Age (Idade):} Valor exato (e.g., 64 anos) $\rightarrow$ Faixas de 5 anos (e.g., 60-65) $\rightarrow$ Faixas de 10 anos (e.g., 60-70).
    \item \textbf{Ethnicity (Etnia):} Categorias específicas (Caucasian, African American, Asian, etc.) $\rightarrow$ Supressão (indicada como "*").
    \item \textbf{Gender (Sexo):} Preservado em seu valor original em todos os cenários (Male / Female).
\end{itemize}

\subsection{Configuração dos Cenários de Privacidade}

Foram definidos três cenários de anonimização progressiva, aplicados à amostra total ($N = 42.277$) antes da filtragem de duração mínima de permanência, visando fornecer máxima diversidade demográfica para as métricas de risco:

\begin{table}[ht]
\centering
\caption{Configuração dos Cenários Experimentais de Anonimização}
\label{tab:cenarios-anonimizacao}
\small
\begin{tabular}{l|p{3cm}|p{2.5cm}|p{2.5cm}|p{2cm}}
\hline
\textbf{Scenario} & \textbf{Technique / Config.} & \textbf{Age Treatment} & \textbf{Ethnicity Treatment} & \textbf{Gender} \\ \hline
\textbf{Baseline} & Original Data (No anonymization) & Exact Value & Original (Detailed) & Original \\ \hline
\textbf{Scenario A} & $k$-Anonymity ($k=5$) \newline Outlier Suppression & 5-year bands & Suppressed (*) & Original \\ \hline
\textbf{Scenario B} & $k$-Anonymity ($k=10$) \newline $l$-Diversity ($l=2$) & 10-year bands & Suppressed (*) & Original \\ \hline
\textbf{Scenario C} & Differential Privacy \newline ($\epsilon=2.0$, $\delta=10^{-5}$) & 10-year bands & Original (Preserved) & Original \\ \hline
\end{tabular}
\end{table}

A justificativa para cada cenário é a seguinte:

\begin{itemize}
    \item \textbf{Scenario A ($k=5$):} Representa uma proteção sintática leve, focando em alta utilidade com risco residual moderado. A supressão de Ethnicity reflete a dificuldade computacional de manter granulosidade em faixas etárias finas (5anos) enquanto se atinge $k=5$. Este cenário testa o balanço recomendado por diretrizes clínicas iniciais de de-identificação \cite{ElEmam2011}.
    
    \item \textbf{Scenario B ($k=10, l=2$):} Representa uma proteção sintática robusta, onde o aumento de $k$ para 10 combinado com l-diversidade fornece garantias semânticas contra ataques de homogeneidade. A faixa de 10 anos para Age é necessária para atender simultaneamente os critérios de $k=10$ e $l=2$. Este cenário alcança um equilíbrio conservador recomendado para dados clínicos sensíveis \cite{Machanavajjhala2007}.
    
    \item \textbf{Scenario C (Differential Privacy, $\epsilon=2.0$):} Representa uma proteção probabilística rigorosa. O valor $\epsilon=2.0$ foi selecionado como ponto de equilíbrio empírico em aplicações de privacidade diferencial em saúde, oferecendo garantia matemática forte ("plausible deniability") enquanto preserva a viabilidade de treinamento de modelos complexos. A preservação de Ethnicity em valores originais reflete a estratégia de mitigar potenciais enviesamentos algorítmicos por meio da precisão demográfica, ao custo da amostra (discussão de data loss abaixo).
\end{itemize}

\section{Avaliação de Risco de Reidentificação e Utilidade Clínica}
\label{sec:metod-avaliacao}

A avaliação foi conduzida em duas dimensões complementares: (1) quantificação do risco de reidentificação residual; e (2) avaliação da utilidade clínica através do desempenho do modelo preditivo. Esta estrutura bidirecional permite responder à questão central de pesquisa sobre a viabilidade de equilibrar privacidade e utilidade em dados clínicos.

\subsection{Auditoria de Risco de Reidentificação}

Cada dataset gerado nos três cenários foi submetido a uma auditoria de risco de reidentificação usando a ferramenta ARX. A auditoria modelou três perfis de ataque distintos, conforme a taxonomia de El Emam \cite{ElEmam2011}:

% TODO: adicionar na fundamentação teórica

\begin{enumerate}
    \item \textbf{Prosecutor Risk (Risco de Promotor):} Representa o cenário em que o adversário já sabe que um indivíduo específico está no dataset e tenta provar isso. É calculado como a proporção máxima de risco, equivalente a $1/k_{\min}$, onde $k_{\min}$ é o tamanho do menor grupo de equivalência.
    
    \item \textbf{Journalist Risk (Risco de Jornalista):} Representa um adversário que não tem certeza se o alvo está no dataset. É mitigado aplicando ajustes estatísticos baseados na proporção esperada da população, resultando frequentemente em valores similares ao Prosecutor Risk para pequenas proporções de amostragem.
    
    \item \textbf{Marketer Risk (Risco de Profissional de Marketing):} Representa um adversário sem alvo específico tentando reidentificar o máximo possível de registros. É medido como a taxa média esperada de sucesso ($E[1/k]$) em toda a população.
\end{enumerate}

Adicionalmente, foram registradas a proporção de registros únicos (registros com $k=1$) antes e depois da anonimização, bem como o impacto em termos de redução de amostra (relevante especialmente para Differential Privacy).

\subsection{Avaliação de Utilidade Clínica}

A utilidade foi avaliada através de uma abordagem orientada à tarefa (\textit{task-based}). O modelo ConCare foi treinado separadamente em cada uma das versões anonimizadas do dataset (Scenarios A, B, C) e nos dados originais (Baseline), utilizando-se a mesma divisão de treinamento/validação/teste.

% TODO: expandir explicação e justificativa para as métricas de avaliação usando literatura e dados reais

Para medir o desempenho preditivo, foram utilizadas as seguintes métricas:

\begin{itemize}
    \item \textbf{AUROC (Area Under the Receiver Operating Characteristic Curve):} Mede a capacidade global do modelo de discriminar pacientes que evoluíram para óbito daqueles que receberam alta, independente de limiar de decisão. Varia de 0 (pior) a 1 (perfeito).
    
    \item \textbf{AUPRC (Area Under the Precision-Recall Curve):} Métrica particularmente rigorosa em contextos de desbalanceamento severo de classes, como é o caso da mortalidade em UTI ($\approx 10.6\%$ na amostra). AUPRC penaliza fortemente falsos positivos, sendo a métrica primária de avaliação.
\end{itemize}

Para garantir a robustez estatística das comparações, foram calculados \textbf{Intervalos de Confiança de 95\%} (IC 95\%) utilizando o método de \textit{bootstrap} não-paramétrico com 1.000 reamostragens sobre as predições do conjunto de teste. Esse procedimento permite avaliar simultaneamente a magnitude do efeito e a sua incerteza, fornecendo uma visão completa da estabilidade das estimativas.

Os resultados de risco e utilidade foram compilados e visualizados em gráficos de trade-off (Risco $\times$ Utilidade) para identificar o ponto ótimo de equilíbrio, sintetizando a viabilidade técnica e jurídica da anonimização para pesquisa clínica sob a LGPD.
